% !TEX root = ../supplement_0421052099.tex
% ---------------------------------------------------------------------
% APPENDIX B. Created 2026-09-23 when the user set a hard budget of
% 20-22 pages for abstract-to-conclusion. The budget excludes the
% appendices, so this is RELOCATED material, not deleted material:
% everything here was in Section 3 and is reachable from it by an
% explicit pointer. Nothing was dropped in the move.
%
% IT LIVES IN sections/ ON PURPOSE. scripts/check_paper_numbers.py
% reads journal/sections/*.tex plus the root file and nothing else, so
% a number moved out of the body into a file under sections/ keeps its
% binding. Moving this to journal/appendix/ would silently break every
% check that pins a constant of the loop. The same applies to
% section_body(), which finds a label wherever it sits.
%
% sec:removal keeps its label even though it is now an appendix
% subsection: check_paper_numbers.py's PROMISES block requires a
% section labelled sec:removal containing the word "flag", and the
% body's pointer to it is \ref (elsarticle's \appendix redefines
% \thesection, so \Cref would print "Section Appendix B.5").
%
% Read off agr/nodes.py, agr/scorer.py, agr/planner.py, agr/budget.py
% and agr/kg_tools.py; where a number is a code constant it is the
% constant the reported runs used.
% ---------------------------------------------------------------------

\section{The instrument in detail}\label{app:instrument}

\Cref{sec:framework} gives the loop at the width the attribution study
needs. This appendix records the mechanics a reimplementation would
want, in the order a question meets them.

\begin{figure}[!tb]
  \begin{center}
    \begin{tikzpicture}[
        node distance=9mm,
        box/.style={draw=agrNode, thick, rounded corners=2pt,
                    minimum height=7.5mm, fill=agrNode!8,
                    minimum width=24mm, align=center, font=\small},
        vbox/.style={box, draw=agrSup, fill=agrSup!12},
        term/.style={draw=black!45, rounded corners=6pt, minimum height=6mm,
                     minimum width=16mm, font=\scriptsize\sffamily,
                     fill=black!6},
        flow/.style={-{Stealth[length=2mm]}, semithick, draw=black!68},
        lbl/.style={font=\scriptsize\ttfamily, inner sep=1.5pt,
                    fill=white, anchor=center}]

      \node[term]              (start) {START};
      \node[box, below=of start] (plan) {Planner};
      \node[box, below=of plan]  (expl) {Explorer};
      \node[box, below=of expl]  (eval) {Evaluator};
      \node[vbox, below=of eval] (ver)  {Verifier};
      \node[box, below=of ver]   (ans)  {Answerer};
      \node[term, below=of ans]  (end)  {END};
      \node[box, right=26mm of eval] (back) {Backtracker};

      \draw[flow] (start) -- (plan);
      \draw[flow] (plan)  -- (expl);
      \draw[flow] (expl)  -- (eval);
      \draw[flow] (eval)  -- node[lbl] {answer} (ver);
      \draw[flow] (ver)   -- node[lbl] {grounded $\mid$ give\_up} (ans);
      \draw[flow] (ans)   -- (end);

      \draw[flow] (eval.west) -- ++(-22mm,0)
        node[lbl, above, pos=0.5] {continue}
        |- (expl.west);

      \draw[flow] (eval.east) -- node[lbl] {backtrack} (back.west);
      \draw[flow] (back.north) |- (expl.east);

      \draw[flow] (ver.west) -- ++(-30mm,0)
        node[lbl, below, pos=0.5] {retry}
        |- ([yshift=-2mm] expl.west);

      \begin{scope}[on background layer]
        \node[draw=black!35, dashed, rounded corners, inner sep=2.5mm,
              fit=(expl)(eval)] {};
      \end{scope}
    \end{tikzpicture}
  \end{center}
  \caption{The AGR state machine. The dashed region is the
    Explorer\,$\leftrightarrow$\,Evaluator search loop.
    Evaluator $\rightarrow$ Backtracker $\rightarrow$ Explorer restores
    an earlier frontier, and Verifier $\rightarrow$ Explorer is
    verification-driven re-exploration. The budgets of
    \ref{sec:budgets} bound all three cycles. Every path from
    Evaluator to Answerer passes through the Verifier.}
  \label{fig:statemachine}
\end{figure}

\subsection{Graph tool API}\label{app:tools}

Four graph operations back the loop, as parameterised Cypher templates
against a Neo4j property graph rather than free-form queries the model
would write itself, because a generated-query failure is
indistinguishable in the results from a reasoning failure.
\textsf{search\_entity} resolves a surface form to a node through three
tiers, exact name, then full-text index, then vector similarity, and
stops at the first that returns a match. \textsf{get\_relations} returns
the relation types incident on an entity with their direction and
fanout, after a blocklist removes schema relations, sorted by descending
fanout and cut at $300$ rows. \textsf{get\_neighbors} returns at most
$200$ neighbours along one relation. Unnamed mediator nodes, which
encode $n$-ary facts and make up $63.7\%$ of the graph
(\Cref{sec:environment}), are traversed transparently, so a mediator
path costs one hop. \textsf{verify\_connection} asks whether two
entities are adjacent directly or through one mediator, ignoring the
relation. A fifth operation, \textsf{verify\_triple}, was built first to
check a claim $\langle h, r, t \rangle$ by asking whether that exact
relation held between the two entities. It did not survive contact with
generated text, because claim decomposition phrases relations in English
(``played for'', ``is the capital of'') and those phrases rarely
correspond to a single Freebase predicate, so matching on them rejected
claims that were true. Nothing in the running loop calls it.

\subsection{Edge scoring}\label{app:scoring}

The explorer gathers every relation row of every anchor into one
candidate set, at most five anchors, and scores it in one pass. Each
row $c$, an anchor together with a relation and a direction, receives
\begin{equation}\label{eq:score}
  \mathrm{score}(c) \;=\; \alpha \cos\!\big(v_{\mathrm{rel}(c)},\,
  v_{o}\big) \;+\; (1-\alpha)\, \mathrm{LLM}(c),
\end{equation}
which blends two signals at the frozen $\alpha = 0.7$. The first is the
cosine similarity between the candidate relation's embedding and the
embedding of the active sub-objective $o$, both from
\textsf{all-MiniLM-L6-v2}~\cite{reimers2019sentence}, with the relation
vocabulary of $7{,}058$ types embedded once in advance. The objective,
not the original question, is what gets embedded, and that substitution
is the concrete mechanism through which decomposition is meant to help.
The second is a batched model judgement, one call, over the twenty
candidates with the highest embedding score; every other candidate
receives zero for that term. The blend is applied after the model call
rather than inside the prompt, so that a prompt free of any
configuration value lets an $\alpha$ sweep reuse one set of cached
responses across every condition. The three best rows per anchor,
skipping banned ones, are expanded, and the five neighbours with the
highest reaching score become the next frontier.

\subsection{Evaluation, routing, and backtracking}\label{app:routing}

The evaluator is one model call. It sees the active sub-objective, the
thirty traversed triples most similar to it by the same embedding, and
the remaining depth and backtrack budget, and it returns one of three
decisions together with the entities it considers to satisfy the
objective. Those entities are accepted only if they are names in the
traversed set, so the evaluator cannot introduce an entity the graph did
not return. When it marks an objective done, its resolved entities
become the anchors for the next objective.

The router after it applies three tests in order. If the depth or the
wall-clock budget is spent, the loop ends and the verifier runs.
Otherwise, if a backtrack trigger fires and the backtrack budget allows,
the backtracker runs; three triggers exist, namely that the last
expansion added no triple, that the best score of the last expansion
fell below $\tau = 0.20$, or that the evaluator asked to backtrack.
Otherwise the evaluator's decision stands. The backtracker pops the
snapshot with the highest recorded score from the stack, which is often
not the most recent one, restores its anchors and depth, and bans the
anchor--relation--direction rows the most recent expansion followed. The
mismatch between the snapshot it restores and the rows it bans is
deliberate in neither direction, and \Cref{sec:nulls} measures what it
costs. This is where the mechanics part from
Plan-on-Graph~\cite{chen2024plan}, whose decomposition and reflection
live in prompts the model executes: here the plan is a data structure
the program advances, two of the three backtrack triggers are
deterministic tests on the scores, and the ban list, not a reflection,
is what makes a restored frontier explore differently.

\subsection{Claim routing and the verifier's outcomes}\label{app:claims}

One model call sees the question, the sixty traversed triples most
similar to it, and up to twenty-five candidates. It drafts prose,
extracts the entities the draft asserts as its answer, and decomposes
the draft into atomic claims over entity \emph{names}, since the model
has never seen a graph identifier; the relation slot is left as free
text for the reason \ref{app:tools} gives.

Each claim is routed through up to three tests and leaves at the first
that settles it. \textbf{Traversed adjacency} collapses the traversed
triples into an undirected set of endpoint pairs; a claim is supported
if its pair is in that set, and every traversed triple joining the pair
is attached as citable evidence. This is the only route of the three
that attaches evidence. \textbf{\textsf{verify\_connection}} catches
claims whose endpoints were both reached but never walked together,
querying the full graph for the same relation-agnostic adjacency. Claims
that clear neither route, together with those naming an entity the agent
never traversed at all, fall to an \textbf{entailment check}: one
batched model call judging each remaining claim against the same fact
window the drafter saw. Symbols decide, and the model is consulted only
where the graph is silent.

The router that follows has three outcomes. No unsupported claim routes
to \textsf{grounded}, and the draft is emitted verbatim with no further
model call. An unsupported claim with verification, depth, and time
budget remaining routes to \textsf{retry}: the verifier appends a repair
objective naming the first unsupported claim, anchors the frontier on
that claim's head entity, and control returns to the Explorer. Otherwise
the router takes \textsf{give\_up}, and one rewrite call restates the
draft around a \emph{deterministically filtered} entity list, in which
an entity survives if it appears in a supported claim or in no
unsupported claim. An earlier design let the model regenerate the list
itself, and that reintroduced exactly the ungroundedness the layer
exists to prevent.

On the $80$-question development set used to validate this design, $77$
questions asserted no unsupported claim and took the grounded path at
zero additional model cost. The remaining $3$ took \textsf{give\_up},
and the retry route did not fire once, because every question that
reached the verifier with an unsupported claim had already exhausted its
depth budget. On most questions, then, the layer is a pass-through that
certifies a draft rather than a mechanism that visibly intervenes, which
is the reading \Cref{sec:nulls} confirms at test scale.

\subsection{Budgets and termination}\label{sec:budgets}

Routers and the model client enforce every cap; no prompt requests one.
A run may take at most $4$ expansion depths, a beam of $3$ candidate
rows per anchor, $5$ anchors in the frontier, $3$ backtracks, $2$
verification iterations, $300$ seconds of wall-clock time, and $25$
model calls. The depth, wall-clock, and verification caps are checked at
two separate sites each, because the \textsf{retry} edge re-enters the
Explorer without passing back through the router that would otherwise
catch it. Every cap is hashed into the run configuration, so a reported
result can always be traced to the caps that produced it. The two
tool-layer caps of \ref{app:tools}, $300$ candidate relations and $200$
neighbours, bound what the agent can see rather than what it can spend,
and \Cref{sec:baselines} shows why they matter for the comparison. Every
cycle in \Cref{fig:statemachine} therefore passes through a router that
checks or decrements a monotone counter, so no sequence of model
outputs, however adversarial, can produce a non-terminating run.

\subsection{Removing a component}\label{sec:removal}

Each ablation condition of \Cref{sec:attribution} flips one flag and
leaves every other line of the loop as it is. \textbf{No planner}
replaces the planning call with a plan of one objective, the question
itself; anchors still come from the annotated topic entities, as they do
in the full system. \textbf{Embedding-only scoring} sets $\alpha = 1$,
which makes no scoring call and ranks candidates by cosine similarity
alone. \textbf{No backtracking} sets the backtrack budget to zero, so
the router ends the search where it would have backtracked. \textbf{No
verification} leaves the drafting call unchanged, entity list and claim
decomposition included, but checks no route, makes no entailment call,
and emits the draft verbatim. That condition therefore measures the
checking and not the drafting, and its token saving is the cost of the
routes rather than of the draft. The budget configuration is hashed into
every record, so the condition that produced a row of
\Cref{tab:ablation} can be recovered from the row.
